\cleardoublepage
\chapter{PENDAHULUAN}
\label{chap:pendahuluan}

\section{Latar Belakang}

Pergerakan harga saham suatu emiten di Bursa Efek Indonesia (BEI) tidak selalu disebabkan oleh berita yang berkaitan langsung dengan perusahaan tersebut. Harga suatu emiten dapat turun karena Indeks Harga Saham Gabungan (IHSG) sedang melemah, karena sektor industrinya sedang tertekan, atau karena terdapat berita yang secara khusus menyangkut perusahaan, seperti kasus hukum, pergantian direksi, atau rilis laporan keuangan. Dalam kajian keuangan, dampak suatu peristiwa terhadap nilai perusahaan umumnya diukur melalui \emph{event study}. Pendekatan ini memanfaatkan asumsi bahwa informasi baru akan tercermin pada harga saham, sehingga dampaknya dapat diukur dari \emph{abnormal return}, yaitu selisih antara \emph{return} aktual dan \emph{return} yang diharapkan berdasarkan kondisi pasar \autocite{mackinlay1997,brown1985}. Bagi investor dan analis, pemisahan antara pengaruh pasar, pengaruh sektor, dan pengaruh berita menjadi penting karena menentukan bagaimana suatu informasi seharusnya ditanggapi.

Selain pengaruh pasar secara umum, pergerakan sektor juga perlu diperhitungkan dalam menghitung \emph{return} yang diharapkan. Sejak 25 Januari 2021, BEI menerapkan klasifikasi IDX Industrial Classification (IDX-IC) yang menggantikan JASICA dan menyediakan sebelas indeks sektoral, seperti IDXENERGY, IDXBASIC, IDXFINANCE, dan IDXTECHNO \autocite{cnbc2021idxic}. Ketersediaan indeks sektoral ini memungkinkan \emph{return} yang diharapkan dihitung tidak hanya dari pergerakan IHSG, tetapi juga dari pergerakan sektor emiten. \textcite{xiang1993} menunjukkan bahwa faktor industri menjelaskan rata-rata 6,7\% variasi residual \emph{market model} pada data harian dan 11,5\% pada data mingguan. Penelitian tersebut juga menyimpulkan bahwa apabila \emph{return} saham sebenarnya mengikuti model dua faktor (pasar dan industri), inferensi yang hanya didasarkan pada \emph{market model} menjadi tidak valid. Dengan demikian, koreksi terhadap pengaruh sektor diperlukan agar bagian \emph{return} yang dikaitkan dengan berita tidak tercampur dengan pergerakan sektor.

Penelitian di Indonesia yang memanfaatkan berita untuk memahami pergerakan saham sejauh ini umumnya berfokus pada analisis sentimen. \textcite{attar2024} menggunakan IndoBERT untuk mengklasifikasikan sentimen 2.857 berita saham dan memperoleh akurasi 0,78, kemudian mengukur korelasinya dengan pergerakan harga per sektor dengan korelasi Pearson tertinggi sebesar 0,389 pada sektor \emph{basic materials}. \textcite{alamsyah2025} menganalisis 192.582 berita keuangan Indonesia periode 2018 sampai 2023 dan menemukan bahwa korelasi sentimen dengan harga meningkat dari 0,26 pada tingkat perusahaan menjadi 0,30 pada tingkat industri dan 0,43 pada klaster emiten yang sering muncul bersama di berita. \textcite{kusuma2026} menggabungkan sentimen IndoBERT dari 9.819 berita CNBC Indonesia dengan data harga untuk mengklasifikasikan tren harga lima emiten LQ45 menggunakan LSTM. Ketiga penelitian tersebut menunjukkan bahwa berita berbahasa Indonesia memiliki keterkaitan dengan pergerakan harga saham. Namun, informasi berita yang digunakan masih terbatas pada sentimen, sedangkan targetnya berupa harga atau tren harga mentah.

Penggunaan sentimen saja memiliki keterbatasan karena sentimen hanya menggambarkan nada berita, bukan jenis peristiwa yang terjadi. Berita mengenai direktur yang ditetapkan sebagai tersangka korupsi dan berita mengenai laba yang turun sesuai ekspektasi analis sama-sama bernada negatif, tetapi keduanya merupakan peristiwa yang berbeda. \textcite{zhou2021} memandang peristiwa korporasi sebagai penggerak pergerakan saham dan membangun model deteksi peristiwa dari berita sebagai dasar strategi \emph{trading}, berbeda dengan model yang langsung memanfaatkan fitur teks dan sentimen. \textcite{dolphin2026} juga melakukan \emph{event study} terhadap besar \emph{abnormal return} tanpa memperhatikan arahnya pada laporan 8-K dan menunjukkan bahwa kategori peristiwa yang lebih rinci dapat membedakan peristiwa korporasi yang secara ekonomi berbeda. Selain itu, penggunaan harga mentah sebagai target membuat pengaruh berita bercampur dengan pengaruh pasar dan sektor. Akibatnya, model dapat mempelajari pola yang sebenarnya berasal dari pergerakan IHSG atau sektor, bukan dari berita.

Penelitian di Indonesia yang sudah menggunakan \emph{abnormal return} juga masih memiliki keterbatasan. \textcite{elli2025} menguji hubungan sentimen 23.406 \emph{tweet} dengan \emph{abnormal return} saham IDX30 pada periode Juli sampai Desember 2023 menggunakan uji kausalitas Granger. Sumber data yang digunakan berupa \emph{tweet} dan analisisnya bersifat statistik, bukan model prediksi untuk setiap berita. Di sisi lain, beberapa \emph{event study} di BEI meneliti satu atau dua jenis peristiwa yang ditentukan secara manual, misalnya \emph{stock split} pada periode 2015 sampai 2019 \autocite{fauzi2021} serta \emph{right issue} dan \emph{stock split} pada periode 2022 sampai 2024 \autocite{sudirman2026}. Pada penelitian \textcite{sudirman2026}, \emph{abnormal return} dihitung sebagai selisih \emph{return} saham terhadap \emph{return} IHSG, sehingga pengaruh sektor belum dikoreksi. Pendekatan tersebut juga belum dapat digunakan untuk menilai berita baru secara otomatis karena peristiwa dan tanggalnya harus ditentukan terlebih dahulu oleh peneliti.

Pada tingkat global, otomatisasi deteksi peristiwa dari teks dan pengukuran \emph{abnormal return} sudah mulai dilakukan. \textcite{kargarzadeh2026} menggunakan \emph{large language model} (LLM) sebagai \emph{teacher} yang kemudian didistilasi menjadi \emph{classifier} untuk memberi 17 label peristiwa pada 4,57 juta artikel berita, lalu mengukur \emph{abnormal return} yang disesuaikan dengan beta pada sekitar 3.000 saham Amerika Serikat. \textcite{dolphin2026} juga memanfaatkan LLM untuk memberi label 119 jenis peristiwa pada sekitar 293.000 laporan 8-K. Kedua penelitian tersebut berfokus pada pasar Amerika Serikat, dan \emph{abnormal return} pada penelitian \textcite{kargarzadeh2026} hanya dikoreksi terhadap pengaruh pasar melalui beta. Penggunaan LLM juga memiliki risiko tersendiri. \textcite{gao2025} menemukan bahwa LLM menunjukkan kecenderungan \emph{look-ahead} yang tinggi pada periode data pelatihannya, tetapi kecenderungan tersebut hampir hilang setelah batas waktu data pelatihan. Hal ini menunjukkan bahwa kinerja prediksi LLM pada data historis dapat dipengaruhi oleh informasi mengenai hasil yang sebenarnya sudah terjadi.

Selain akurasi, keluaran model prediksi juga perlu dapat dijelaskan agar pengguna memahami faktor yang mendasari suatu prediksi. SHAP (\emph{SHapley Additive exPlanations}) memberikan nilai kontribusi setiap fitur terhadap suatu prediksi tertentu \autocite{lundberg2017}, sehingga dapat digunakan untuk menjelaskan keluaran model berbasis pohon seperti XGBoost \autocite{chen2016}. \textcite{ruan2025} telah menggabungkan sentimen FinBERT dengan indikator teknikal pada model XGBoost dan menggunakan SHAP untuk menjelaskan prediksi pergerakan harga jangka pendek saham S\&P 500. Namun, fitur berita yang digunakan masih berupa sentimen dan targetnya masih berupa pergerakan harga, sehingga SHAP belum dapat memisahkan kontribusi berita dari kontribusi pasar. Di sisi lain, prediksi arah pergerakan harga dari berita memiliki tingkat kesulitan yang tinggi. \textcite{luss2015} menemukan bahwa pada prediksi \emph{return} intrahari dari berita, arah \emph{return} tidak dapat diprediksi menggunakan teks maupun \emph{return} historis, tetapi besarnya pergerakan dapat diprediksi dan fitur teks memberikan kinerja yang lebih baik dibandingkan \emph{return} historis saja. Oleh karena itu, model perlu dievaluasi tidak hanya dari ketepatan arah, tetapi juga dari estimasi besarnya \emph{abnormal return}.

Berdasarkan kondisi tersebut, penelitian ini mengusulkan sistem yang menghubungkan berita, jenis peristiwa, emiten, dan pergerakan harga dalam satu \emph{pipeline}. Sistem mendeteksi emiten yang disebutkan dalam berita menggunakan kamus kode, nama, dan alias emiten, kemudian mengklasifikasikan jenis peristiwa menggunakan IndoBERT yang di-\emph{fine-tune} \autocite{wilie2020}. Fitur berita, seperti jenis peristiwa, sentimen, keterlibatan emiten, intensitas, dan waktu terbit, digabungkan dengan fitur pasar yang tersedia sebelum berita terbit. Target yang diprediksi adalah \emph{abnormal return} yang dikoreksi terhadap IHSG dan indeks sektor, sedangkan model yang digunakan adalah XGBoost untuk memprediksi arah dan besar \emph{abnormal return}, dengan SHAP sebagai penjelasan setiap prediksi. Seluruh \emph{pipeline} dirancang tanpa LLM generatif untuk menjaga \emph{reproducibility}, menghindari risiko \emph{look-ahead bias}, dan menjaga transparansi rantai penjelasan. Berdasarkan penelusuran literatur yang telah dilakukan, belum ditemukan penelitian yang secara bersamaan mendeteksi jenis peristiwa dari berita berbahasa Indonesia secara otomatis, mengestimasi \emph{abnormal return} emiten BEI setelah mengendalikan IHSG dan indeks sektor, serta menjelaskan kontribusi setiap faktor pada tingkat berita.

\section{Rumusan Masalah}

Penelitian yang memanfaatkan berita saham berbahasa Indonesia masih menggunakan sentimen sebagai satu-satunya informasi dari berita dan harga mentah sebagai target, sehingga dampak berita tidak dapat dipisahkan dari pengaruh pasar dan sektor. Sementara itu, pengukuran \emph{abnormal return} di BEI masih dilakukan melalui \emph{event study} dengan peristiwa yang ditentukan secara manual. Akibatnya, model yang dilatih dengan target harga mentah dapat mempelajari pola yang berasal dari pergerakan IHSG atau sektor, bukan dari berita, sehingga hasilnya dapat menyesatkan ketika digunakan untuk menilai dampak suatu berita. Berdasarkan penelusuran literatur yang telah dilakukan, belum ditemukan sistem yang memberikan estimasi terukur mengenai keterkaitan suatu berita berbahasa Indonesia dengan pergerakan emiten BEI setelah pengaruh IHSG dan sektor dikendalikan. Penelitian ini mengusulkan sistem berbasis klasifikasi jenis peristiwa, target \emph{abnormal return} dua faktor, dan penjelasan SHAP untuk menjawab permasalahan tersebut. Berdasarkan uraian tersebut, rumusan masalah dalam penelitian ini adalah sebagai berikut.
\begin{enumerate}
    \item Sejauh mana jenis peristiwa yang diekstraksi secara otomatis dari berita berbahasa Indonesia memberikan informasi tambahan dibandingkan sentimen dalam memprediksi arah dan besar \emph{abnormal return} emiten BEI?
    \item Bagaimana perbedaan kinerja model dan pola kontribusi fitur yang dihasilkan SHAP ketika target yang digunakan adalah \emph{abnormal return} yang dikoreksi IHSG dan indeks sektor dibandingkan dengan \emph{return} mentah?
    \item Seberapa baik \emph{classifier} jenis peristiwa berbasis IndoBERT dibandingkan dengan aturan kata kunci serta TF-IDF dengan Logistic Regression atau SVM dalam mengklasifikasikan jenis peristiwa pada berita saham berbahasa Indonesia?
    \item Faktor berita dan kondisi pasar apa saja yang memberikan kontribusi terbesar terhadap estimasi \emph{abnormal return} berdasarkan analisis SHAP?
\end{enumerate}

\section{Tujuan}

Tujuan penelitian ini adalah sebagai berikut.
\begin{enumerate}
    \item Membangun \emph{pipeline} yang menerima teks atau URL berita dan menghasilkan emiten terdampak, jenis peristiwa, arah \emph{abnormal return} beserta probabilitasnya, estimasi besar \emph{abnormal return}, kandidat emiten yang berpotensi terkena \emph{spillover}, serta penjelasan faktor penyumbang prediksi.
    \item Menyusun dataset berita saham berbahasa Indonesia yang berlabel jenis peristiwa, termasuk \emph{test set} emas yang dilabeli secara manual oleh dua anotator.
    \item Membandingkan kinerja aturan kata kunci, TF-IDF dengan Logistic Regression atau SVM, dan IndoBERT dalam mengklasifikasikan jenis peristiwa.
    \item Mengukur kontribusi jenis peristiwa terhadap prediksi \emph{abnormal return} melalui \emph{ablation study} yang membandingkan model dengan fitur sentimen saja, model dengan sentimen dan jenis peristiwa, serta model dengan seluruh fitur.
    \item Membandingkan hasil pemodelan dengan target \emph{abnormal return} dan target \emph{return} mentah dari sisi kinerja prediksi dan pola kontribusi fitur.
    \item Menyajikan penjelasan prediksi menggunakan SHAP, baik secara global maupun untuk setiap berita.
\end{enumerate}

Keberhasilan penelitian ditentukan oleh beberapa kriteria. Pertama, \emph{pipeline} dapat dijalankan secara utuh dari input berita hingga keluaran prediksi dan penjelasan SHAP melalui prototipe antarmuka web. Kedua, kualitas anotasi \emph{test set} emas mencapai nilai Cohen's kappa sekurang-kurangnya 0,6. Ketiga, seluruh model dievaluasi pada data uji periode 2025 yang tidak digunakan dalam pelatihan maupun penyetelan, dan dibandingkan dengan model \emph{baseline} menggunakan metrik yang telah ditetapkan. Keempat, \emph{ablation study} dan perbandingan target dapat dilaporkan beserta uji signifikansi selisih kinerjanya. Hasil yang menunjukkan bahwa jenis peristiwa tidak menambah kinerja model tetap dianggap sebagai temuan yang sah dan dilaporkan apa adanya.

\section{Ruang Lingkup, Asumsi, dan Batasan Masalah}

Ruang lingkup penelitian ini mencakup emiten konstituen LQ45 atau IDX80 yang akan ditetapkan setelah audit ketersediaan data, dengan mengikuti daftar konstituen pada setiap periode untuk menghindari \emph{survivorship bias}. Periode data yang digunakan adalah 2021 sampai 2025, menyesuaikan dengan ketersediaan indeks sektoral IDX-IC sejak Januari 2021. Karena estimasi parameter \emph{abnormal return} membutuhkan data historis sebelum berita, sampel berita direncanakan mulai sekitar Juli 2021. Data berita dikumpulkan dari dua sampai tiga portal berita keuangan nasional, sedangkan data harga saham, IHSG, dan indeks sektoral diperoleh dari sumber publik. Jenis peristiwa dibatasi pada enam kategori utama, yaitu hukum atau korupsi, laporan keuangan, pergantian direksi, kebijakan pemerintah, aksi korporasi, dan operasional, ditambah kategori pergerakan harga dan kategori lainnya.

Penelitian ini menggunakan beberapa asumsi. \emph{Abnormal return} diasumsikan sebagai proksi dampak berita, yaitu bagian dari \emph{return} emiten yang tidak dapat dijelaskan oleh pergerakan IHSG dan indeks sektornya. Waktu terbit berita diasumsikan tercatat dengan akurat hingga satuan menit, sehingga berita yang terbit setelah perdagangan ditutup sekitar pukul 16.00 WIB atau pada hari libur dapat dialokasikan ke hari bursa berikutnya. Data harga yang digunakan diasumsikan telah disesuaikan terhadap \emph{stock split} dan dividen. Parameter \emph{return} yang diharapkan diasumsikan dapat diestimasi dari data historis sebelum berita terbit.

Penelitian ini memiliki beberapa batasan. Sistem tidak dimaksudkan untuk menyatakan bahwa suatu berita secara mutlak menyebabkan perubahan harga, melainkan memberikan estimasi keterkaitan statistik berdasarkan pola historis. Berita yang hanya melaporkan pergerakan harga dikeluarkan dari data pelatihan model \emph{abnormal return} untuk mencegah \emph{data leakage}. Analisis \emph{spillover} dibatasi pada relasi sektor, grup konglomerasi, dan kemunculan bersama di berita karena data relasi rantai pasok antaremiten di Indonesia terbatas. Untuk peristiwa berskala sektor, seperti kebijakan pemerintah yang berlaku pada satu sektor, \emph{abnormal return} juga dihitung dengan koreksi IHSG saja karena koreksi sektor dapat menghapus dampak peristiwa tersebut. Sampel yang jatuh pada hari suspensi tidak digunakan, dan keterbatasan akibat mekanisme auto rejection juga dicatat sebagai batasan. Input berupa URL hanya didukung untuk portal yang sama dengan sumber dataset, dan seluruh \emph{pipeline} tidak menggunakan LLM generatif.

\section{Hipotesis Tugas Akhir}

Hipotesis penelitian ini disusun berdasarkan beberapa premis. Premis pertama, jenis peristiwa yang berbeda dapat diikuti oleh besar \emph{abnormal return} yang berbeda \autocite{dolphin2026}. Premis kedua, sentimen hanya menggambarkan nada berita sehingga berita dengan sentimen yang sama dapat berasal dari peristiwa yang berbeda. Premis ketiga, faktor industri menjelaskan sebagian variasi \emph{return} yang tidak tertangkap oleh \emph{market model} \autocite{xiang1993}. Premis keempat, arah \emph{return} sulit diprediksi dari teks, sedangkan besarnya pergerakan lebih dapat diprediksi \autocite{luss2015}. Dengan premis tersebut, informasi jenis peristiwa diharapkan dapat menambah kemampuan model dalam memprediksi \emph{abnormal return}, walaupun kinerja prediksi arah yang moderat tetap wajar dalam domain ini.

Hipotesis penelitian ini adalah sebagai berikut.
\begin{quote}
\textbf{H1a:} Model yang menggunakan fitur sentimen dan jenis peristiwa menghasilkan F1-macro dan AUC yang lebih tinggi secara signifikan dibandingkan model yang hanya menggunakan fitur sentimen dalam memprediksi arah \emph{abnormal return} pada data uji periode 2025. Signifikansi selisih diuji menggunakan uji McNemar atau \emph{bootstrap}.
\end{quote}
\begin{quote}
\textbf{H1b:} Model yang menggunakan fitur sentimen dan jenis peristiwa menghasilkan MAE yang lebih rendah secara signifikan dibandingkan model yang hanya menggunakan fitur sentimen dalam memprediksi besar \emph{abnormal return} pada data uji periode 2025. Signifikansi selisih diuji menggunakan \emph{bootstrap}.
\end{quote}

Hipotesis tambahan yang diuji adalah sebagai berikut.
\begin{quote}
\textbf{H2:} \emph{Classifier} IndoBERT yang melalui tahap pemanasan dengan dataset topik keuangan menghasilkan F1-macro yang lebih tinggi dibandingkan aturan kata kunci dalam mengklasifikasikan jenis peristiwa pada \emph{test set} emas.
\end{quote}

\section{Metodologi}

Tahapan penelitian yang dilakukan adalah sebagai berikut.
\begin{enumerate}
    \item \textbf{Investigasi masalah dan studi literatur.} Mengidentifikasi permasalahan serta mempelajari penelitian mengenai \emph{event study}, \emph{abnormal return}, analisis berita saham di Indonesia, deteksi peristiwa dari teks, IndoBERT, \emph{gradient boosting}, dan SHAP. Penelusuran literatur dilakukan melalui mesin pencari web, OpenAlex, dan Garuda dengan kata kunci seperti \emph{abnormal return} berita, IndoBERT saham, dan \emph{news event abnormal return} Indonesia. Penelusuran melalui Google Scholar dan Digilib ITB akan dilengkapi secara manual.
    \item \textbf{Analisis masalah dan kebutuhan.} Menganalisis kondisi penelitian saat ini berdasarkan hasil studi literatur, mengaudit ketersediaan dan kualitas data berita serta data harga, dan menentukan kebutuhan sistem, kategori peristiwa, serta kriteria keberhasilan.
    \item \textbf{Perancangan solusi.} Merancang arsitektur \emph{pipeline} secara bertahap dari alur umum hingga komponen yang lebih rinci, meliputi alur data, pelabelan, perhitungan \emph{abnormal return}, fitur, model, eksperimen, dan prototipe, menggunakan diagram alur dan diagram komponen.
    \item \textbf{Pengumpulan data.} Menyusun tabel referensi emiten, alias, sektor IDX-IC, dan konstituen indeks per periode, kemudian mengumpulkan berita dari portal terpilih dan mengunduh data harga yang telah disesuaikan untuk emiten, IHSG, dan indeks sektoral.
    \item \textbf{Prapemrosesan data.} Membersihkan teks berita, menghapus duplikat, mendeteksi emiten dengan kamus alias, menentukan hari peristiwa (H0) berdasarkan waktu terbit, serta membentuk sampel berupa pasangan berita dan emiten.
    \item \textbf{Pelabelan jenis peristiwa.} Menyusun pedoman anotasi, membuat \emph{test set} emas sebanyak 300 sampai 500 berita yang dilabeli dua anotator dengan target Cohen's kappa sekurang-kurangnya 0,6, serta menyusun data latih sebanyak 500 sampai 800 berita dari pra-label aturan kata kunci yang diperiksa dan dikoreksi secara manual.
    \item \textbf{Perhitungan \emph{abnormal return}.} Mengestimasi parameter $\alpha$, $\beta_1$, dan $\beta_2$ dengan regresi OLS pada jendela estimasi H-130 sampai H-11, kemudian menghitung \emph{abnormal return} pada H0, \emph{cumulative abnormal return} pada H0 sampai H+1, dan arah \emph{abnormal return} berdasarkan ambang standar deviasi residual.
    \item \textbf{Pengembangan model.} Melakukan \emph{fine-tuning} IndoBERT secara dua tahap untuk klasifikasi jenis peristiwa, mengekstraksi fitur berita dan fitur pasar yang hanya menggunakan data sampai H-1, serta melatih XGBoost \emph{classifier} untuk arah dan XGBoost \emph{regressor} untuk besar \emph{abnormal return}. Data dibagi secara kronologis, yaitu Juli 2021 sampai 2023 untuk pelatihan, 2024 untuk validasi, dan 2025 untuk pengujian.
    \item \textbf{Eksperimen dan evaluasi.} Melakukan perbandingan metode klasifikasi jenis peristiwa, \emph{ablation study} fitur, perbandingan target \emph{abnormal return} dengan \emph{return} mentah, serta perbandingan dengan model \emph{baseline}. Evaluasi menggunakan accuracy dan F1-macro untuk klasifikasi jenis peristiwa, F1-macro dan AUC untuk arah, MAE dan RMSE untuk besar \emph{abnormal return}, serta \emph{bootstrap} atau uji McNemar untuk signifikansi selisih kinerja.
    \item \textbf{Analisis interpretabilitas.} Menganalisis nilai SHAP secara global dan pada studi kasus per berita, kemudian memeriksa kesesuaiannya dengan intuisi domain.
    \item \textbf{Pengembangan prototipe.} Membangun antarmuka web sederhana yang menerima teks atau URL berita dan menampilkan hasil prediksi beserta penjelasannya. Implementasi menggunakan bahasa Python dengan pustaka trafilatura untuk ekstraksi artikel, yfinance untuk data harga, statsmodels untuk regresi, model indobenchmark/indobert-base-p1 dari Hugging Face yang dilatih di Google Colab, XGBoost, SHAP, dan Streamlit untuk antarmuka.
    \item \textbf{Penarikan kesimpulan.} Menyusun kesimpulan berdasarkan hasil evaluasi dan pengujian hipotesis, serta memberikan saran untuk penelitian selanjutnya.
\end{enumerate}

\section{Sistematika Penulisan}

Sistematika penulisan laporan tugas akhir ini adalah sebagai berikut.
\begin{enumerate}
    \item \textbf{Bab I Pendahuluan} berisi latar belakang, rumusan masalah, tujuan, ruang lingkup, asumsi, batasan masalah, hipotesis tugas akhir, metodologi, dan sistematika penulisan.
    \item \textbf{Bab II Studi Literatur} berisi pembahasan mengenai \emph{event study} dan \emph{abnormal return}, model \emph{return} yang diharapkan, analisis sentimen dan deteksi peristiwa dari berita, IndoBERT, \emph{gradient boosting}, SHAP, serta penelitian terdahulu yang relevan.
    \item \textbf{Bab III Analisis Masalah} berisi analisis kondisi penelitian saat ini, analisis kebutuhan sistem, analisis sumber data, serta analisis alternatif pendekatan, termasuk perbandingan pendekatan dengan dan tanpa LLM.
    \item \textbf{Bab IV Desain Konsep Solusi} berisi rancangan \emph{pipeline}, rancangan pelabelan, rancangan perhitungan \emph{abnormal return}, rancangan fitur dan model, rancangan eksperimen, serta rancangan prototipe.
    \item \textbf{Bab V Rencana Selanjutnya} berisi rencana pelaksanaan, jadwal kegiatan, kebutuhan yang perlu disiapkan, serta risiko dan mitigasinya.
    \item \textbf{Daftar Pustaka} berisi seluruh referensi yang digunakan dalam penyusunan laporan tugas akhir.
    \item \textbf{Lampiran} berisi dokumen pendukung, seperti pedoman anotasi, struktur dataset, konfigurasi eksperimen, dan hasil evaluasi tambahan.
\end{enumerate}
